\documentclass[11pt,a4paper]{article}

% Fonts: TeX Gyre Pagella with matching math
\usepackage{fontspec}
\usepackage{unicode-math}
\setmainfont{texgyrepagella}[
  Extension      = .otf,
  UprightFont    = *-regular,
  BoldFont       = *-bold,
  ItalicFont     = *-italic,
  BoldItalicFont = *-bolditalic,
]
\setmathfont{texgyrepagella-math.otf}
\frenchspacing

% Page layout
\usepackage[margin=1in]{geometry}

% Math
\usepackage{amsmath,amsthm}

% Tables
\usepackage{booktabs}

% Float placement
\usepackage{float}

% Captions
\usepackage{caption}
\captionsetup{font=small,labelfont=bf}

% Lists
\usepackage{enumitem}

% Prevent orphaned headings
\usepackage{needspace}

% Paragraph formatting
\usepackage{parskip}
\emergencystretch=1em
\hyphenpenalty=200

% Typographic refinement
\usepackage{microtype}

% Tighter bibliography spacing
\usepackage{etoolbox}
\AtBeginEnvironment{thebibliography}{\setlength{\itemsep}{4pt plus 0.5pt}\setlength{\parsep}{0pt}}

% Colours
\usepackage{xcolor}

% Hyperlinks
\usepackage{hyperref}
\hypersetup{
  colorlinks=true,
  allcolors=blue!50!black,
  pdfauthor={K.R. Ryan},
  pdftitle={Parasitic Liquidity}
}

% Running header
\usepackage{fancyhdr}
\pagestyle{fancy}
\fancyhf{}
\fancyhead[L]{\small\itshape K.R.\ Ryan}
\fancyhead[R]{\small\itshape Parasitic Liquidity}
\fancyfoot[C]{\thepage}
\renewcommand{\headrulewidth}{0.4pt}
\fancypagestyle{plain}{%
  \fancyhf{}%
  \fancyfoot[C]{\thepage}%
  \renewcommand{\headrulewidth}{0pt}%
}

% Theorem environments
\theoremstyle{definition}
\newtheorem{definition}{Definition}

% Paragraph spacing
\setlength{\parindent}{0pt}
\setlength{\parskip}{6pt plus 2pt minus 1pt}

\begin{document}
\thispagestyle{plain}

\begin{center}
{\LARGE\bfseries Parasitic Liquidity: Emission Extraction via\\[4pt]
Non-Functional Concentrated Liquidity Positions}

\vspace{12pt}
{\large K.R.\ Ryan}\\[4pt]
{\small Independent Researcher}\\[2pt]
{\small \href{mailto:gancor@gancor.xyz}{gancor@gancor.xyz} \quad|\quad ORCID: 0009-0004-6295-7040}

\vspace{8pt}
April 2026

\vspace{6pt}
{\small\textbf{Keywords:} parasitic liquidity, concentrated liquidity, emission extraction,\\
liquidity mining, gauge design, DeFi security, mechanism design}
\end{center}

\vspace{8pt}
\hrule
\vspace{12pt}

\begin{abstract}
This paper identifies and characterises parasitic liquidity, a class of value extraction from concentrated liquidity decentralised exchanges that distribute token emissions to liquidity providers. Emission gauges reward instantaneous staked liquidity. Swap facilitation requires sustained depth. Parasitic liquidity exploits the gap between these two properties. The mechanism operates at the emission layer and is distinct from sandwich attacks, front-running, just-in-time liquidity, and arbitrage.

The vulnerability is derived from the public source code of two architecturally distinct emission systems (Slipstream CL gauges and PancakeSwap's \texttt{PancakeV3LmPool}) and verified with Foundry proof-of-concept suites against both implementations on Base mainnet. Both inherit the Synthetix reward accumulator pattern, originally designed for fungible token staking, and apply it unchanged to non-fungible concentrated liquidity positions. A position staked for a single block earns the same reward rate per unit of liquidity as a position staked for any longer duration. No warmup, width minimum, or swap-utilisation condition exists in the reward accrual path of either implementation.

Aggregate on-chain analysis reveals that single-tick-width positions rebalancing every block account for a material share of staked liquidity in major pools, reaching 63\% of in-range staked liquidity in one observed WETH/USDC pool.

The paper formalises the profitability conditions, quantifies the dilution imposed on legitimate liquidity providers, and proposes two gauge-level remediation mechanisms verified by fork testing. The vulnerability is present in any CL DEX that distributes emissions proportional to instantaneous staked liquidity without duration, width, or utilisation checks.
\end{abstract}

% ---- 1. Introduction ----
\section{Introduction}\label{sec:intro}

\subsection{Emission distribution in concentrated liquidity DEXs}\label{sec:emission}

Uniswap V3 introduced concentrated liquidity, allowing LPs to specify a price range rather than providing liquidity across the full curve~\cite{adams2021}. Only liquidity spanning the current tick is active for swaps. When protocols added emission incentives to this model, they needed a distribution mechanism. The solution adopted across ve(3,3) DEXs, PancakeSwap, and others was to distribute protocol-native tokens to staked positions in proportion to their active liquidity.

This creates a measurement gap. Traders need swap depth that persists across price movements and absorbs order flow. Gauges reward instantaneous staked liquidity at a single point in time. These are different quantities, and they diverge when positions are narrow, short-lived, or both.

The emission formula in a typical CL gauge is:
\begin{equation}
\text{reward}_i = \text{rewardRate} \times \Delta t \times \frac{L_i}{\sum_j L_j}
\end{equation}
where $L_i$ is position $i$'s liquidity at the active tick and the sum runs over all staked positions spanning that tick. This formula has three properties that together enable parasitic liquidity:

\begin{enumerate}[leftmargin=2em]
\item \textbf{Instantaneous measurement.} $L_i$ is a point-in-time value with no time-averaging or duration component.
\item \textbf{No width requirement.} A position spanning the minimum tick spacing earns the same rate per unit of $L$ as a position spanning an arbitrary number of ticks.
\item \textbf{No utilisation coupling.} Emissions accrue regardless of whether the position facilitated any swaps.
\end{enumerate}

\subsection{Existing extraction taxonomy}\label{sec:taxonomy}

Documented MEV extraction classes all operate at the swap execution layer:

\begin{table}[H]
\centering
\begin{tabular}{@{}lll@{}}
\toprule
Class & Mechanism & Layer \\
\midrule
Sandwich attack & Front-run and back-run a pending swap & Mempool / swap \\
Front-running & Execute before a known pending transaction & Mempool / swap \\
Back-running & Execute immediately after a state change & Post-swap \\
JIT liquidity & Add liquidity before a large swap, remove after & Swap execution \\
Arbitrage & Correct cross-venue price discrepancies & Cross-venue swap \\
\bottomrule
\end{tabular}
\end{table}

Parasitic liquidity operates at the emission distribution layer. It does not require, depend on, or interact with any swap.

\subsection{Contributions}\label{sec:contributions}

This paper makes four contributions:

\begin{enumerate}[leftmargin=2em]
\item Identification of parasitic liquidity as a structurally distinct class of value extraction, with analysis of its origins in the Synthetix reward accumulator (\S\ref{sec:definition}).
\item Formal characterisation of the profitability conditions (\S\ref{sec:analysis}).
\item Foundry proof-of-concept suites against two independent protocol implementations on Base mainnet, with aggregate on-chain evidence of exploitation at scale (\S\ref{sec:poc}).
\item Two remediation mechanisms with analysis of effectiveness and trade-offs (\S\ref{sec:remediation}).
\end{enumerate}

% ---- 2. Parasitic Liquidity ----
\section{Parasitic liquidity}\label{sec:definition}

\subsection{Definition}\label{sec:def}

\begin{definition}[Parasitic Liquidity]
A concentrated liquidity position is parasitic if it satisfies the gauge's staking requirements without providing tradeable depth. The position is withdrawn and redeployed so frequently that it facilitates no swaps, yet it accrues emissions continuously.
\end{definition}

Let $\delta t$ denote the position's hold time per cycle and $\lambda$ the swap arrival rate on the target pool. A position is parasitic when:
\begin{equation}
\lambda \cdot \delta t \ll 1
\end{equation}

\subsection{Strategy}\label{sec:strategy}

A parasitic operator deploys a contract that executes the following loop each block:


\begin{center}\begin{minipage}{0.85\textwidth}\begin{verbatim}
1. Read slot0() on target pool to obtain the current active tick t_c
2. Withdraw existing NFT from gauge
3. Decrease liquidity to zero, collect remaining tokens, burn NFT
4. Mint new NFT with tickLower = t_c, tickUpper = t_c + tickSpacing
5. Deposit new NFT into gauge
6. Repeat next block
\end{verbatim}\end{minipage}\end{center}


On L2 chains with sub-second to two-second block times, this cycle completes every block. The position is always at the current tick, always minimum width, and always staked. Capital recycles each block with no lock-up or impermanent loss.

\subsection{Root cause}\label{sec:rootcause}

The gauge reward formula (Equation 1) is memoryless. A position staked in block $B$ immediately contributes to $\sum_j L_j$ and accrues its proportional share of emissions. When withdrawn in block $B+1$, it collects accrued rewards. The gauge retains no record of staking duration.

\textbf{Synthetix} introduced the reward accumulator pattern in \texttt{StakingRewards.sol} (2019) for fungible token staking. The core mechanism computes a global reward-per-token accumulator:


\begin{center}\begin{minipage}{0.85\textwidth}\begin{verbatim}
function rewardPerToken() public view returns (uint256) {
    if (_totalSupply == 0) return rewardPerTokenStored;
    return rewardPerTokenStored.add(
        lastTimeRewardApplicable().sub(lastUpdateTime)
            .mul(rewardRate).mul(1e18).div(_totalSupply)
    );
}
\end{verbatim}\end{minipage}\end{center}


Each user's earned reward is the delta between the global accumulator and their last-snapshotted value, multiplied by their stake. This pattern was designed for fungible ERC-20 staking where all participants provide an identical product.

\textbf{Curve} adopted this accumulator for gauge-weighted emissions (2020), distributing CRV to LPs in proportion to their pool shares. \textbf{Solidly} (Feb~2022) forked Curve's gauges into the ve(3,3) model with permissionless gauge creation and fees-to-voters. A family of protocols refined this model: \textbf{Velodrome} (Optimism, June~2022), \textbf{Thena} (BNB~Chain, Jan~2023), \textbf{Aerodrome} (Base, August~2023), \textbf{Ramses} (Arbitrum), \textbf{Pharaoh} (Avalanche), \textbf{Shadow} (Sonic), \textbf{SwapX} (Sonic), and others.

The earlier protocols in this lineage operated on constant-product (V2) pools with fungible LP tokens, where a user's LP token represented a proportional claim to the entire pool's liquidity curve. In that context, the instantaneous accumulator was appropriate: all LPs in a pool provided identical liquidity shapes, so rewarding instantaneous stake was equivalent to rewarding sustained depth. Later entrants adopted concentrated liquidity from launch, inheriting the same accumulator without revisiting its assumptions.

The design broke when these protocols adopted concentrated liquidity, grafting the same accumulator onto Uniswap~V3 position mechanics. In~V3, positions are non-fungible and heterogeneous. Two positions with identical liquidity values can provide vastly different tradeable depth depending on their tick range width. The gauge formula rewards both at the same rate per unit of liquidity, regardless of range, duration, or swap utilisation.

PancakeSwap's \texttt{MasterChefV3} arrived at the same vulnerability via parallel evolution, applying the Synthetix accumulator (used since MasterChef~V1) to their Uni~V3 fork. Any protocol that distributes emissions proportional to instantaneous staked liquidity in a concentrated liquidity pool is affected.

The relevant code pattern in Slipstream-derived gauges:


\begin{center}\begin{minipage}{0.85\textwidth}\begin{verbatim}
function _updateRewardsGrowthGlobal() internal {
    uint32 timestamp = _blockTimestamp();
    uint256 timeDelta = timestamp - lastUpdated;
    if (timeDelta != 0) {
        if (rewardReserve > 0) {
            uint256 reward = rewardRate * timeDelta;
            if (reward > rewardReserve) reward = rewardReserve;
            rewardReserve -= reward;
            if (stakedLiquidity > 0) {
                rewardGrowthGlobalX128 += FullMath.mulDiv(
                    reward, FixedPoint128.Q128, stakedLiquidity
                );
            }
        }
        lastUpdated = timestamp;
    }
}
\end{verbatim}\end{minipage}\end{center}


\texttt{stakedLiquidity} is the total active liquidity across all staked positions at the current tick. The equivalent function in PancakeSwap's \texttt{PancakeV3LmPool} uses \texttt{lmLiquidity} for the same purpose:


\begin{center}\begin{minipage}{0.85\textwidth}\begin{verbatim}
function accumulateReward(uint32 currTimestamp) external override {
    if (currTimestamp <= lastRewardTimestamp) return;
    if (lmLiquidity != 0) {
        (uint256 rewardPerSecond, uint256 endTime) =
            masterChef.getLatestPeriodInfo(address(pool));
        uint32 duration;
        if (uint32(endTime) > currTimestamp) {
            duration = currTimestamp - lastRewardTimestamp;
        } else if (uint32(endTime) > lastRewardTimestamp) {
            duration = uint32(endTime) - lastRewardTimestamp;
        }
        if (duration != 0) {
            rewardGrowthGlobalX128 += FullMath.mulDiv(
                duration,
                FullMath.mulDiv(
                    rewardPerSecond, Q128, REWARD_PRECISION
                ),
                lmLiquidity
            );
        }
    }
    lastRewardTimestamp = currTimestamp;
}
\end{verbatim}\end{minipage}\end{center}


% ---- 3. Formal Analysis ----
\section{Formal analysis}\label{sec:analysis}

\subsection{Profitability condition}\label{sec:profit}

Let $E$ be the emission rate per second per unit of active liquidity in USD, $L_p$ the parasitic position's liquidity, $L_{\text{total}}$ the total active liquidity at the current tick, $f$ the rebalance frequency, and $g$ the gas cost per cycle.

Revenue per second:
\begin{equation}
R = E \cdot \frac{L_p}{L_{\text{total}}}
\end{equation}

Cost per second:
\begin{equation}
C = f \cdot g
\end{equation}

Profitability requires:
\begin{equation}
E \cdot \frac{L_p}{L_{\text{total}}} > f \cdot g
\end{equation}

On Base (two-second blocks, sub-cent gas), with $g \approx \$0.005$ per rebalance cycle and one rebalance per block:
\begin{equation}
C = \frac{0.005}{2} = \$0.0025/\text{sec} = \$216/\text{day}
\end{equation}

An operator targeting multiple pools simultaneously amortises the fixed gas cost across all of them. As the capital efficiency analysis in \S\ref{sec:capeff} shows, a small amount of capital can dominate the liquidity share at the active tick.

\subsection{Capital efficiency}\label{sec:capeff}

Each rebalance cycle withdraws and redeploys the full principal, so the operator's capital is never locked. The effective capital requirement is a single position plus gas reserves. Legitimate LPs, by contrast, keep capital deployed and exposed to impermanent loss for the duration of their position.

\subsection{Fee income versus emission income}\label{sec:feevemission}

A single-tick position minted at the current tick covers one tick spacing for one block. Even if swap volume passes through that tick, the fee capture from a two-second window is negligible relative to the emission income earned in the same period. The Foundry tests in \S\ref{sec:poc} verify this empirically.

\subsection{Dilution of legitimate liquidity providers}\label{sec:dilution}

Legitimate LPs receive a reduced fraction of emissions:
\begin{equation}
\frac{L_{\text{legit}}}{L_{\text{legit}} + L_p}
\end{equation}

Single-tick-width positions concentrate all capital into one tick spacing. For small tick ranges, liquidity scales inversely with range width, so the same capital in one tick spacing produces roughly ten times the liquidity of the same capital spread across ten spacings. A small amount of capital can therefore dominate the liquidity share at the active tick and capture a disproportionate fraction of emissions.

\subsection{Scalability}\label{sec:scale}

The parasitic share of emissions scales linearly with the number of operators:
\begin{equation}
\text{Parasitic share} = \frac{n \cdot L_{p,\text{avg}}}{L_{\text{legit}} + n \cdot L_{p,\text{avg}}}
\end{equation}

Nothing in the protocol prevents multiple operators from targeting the same pool. As $n$ increases, legitimate LP share approaches zero. The low capital requirement (\S\ref{sec:capeff}) and sub-cent gas costs on L2s make the barrier to entry negligible.

% ---- 4. Proof-of-Concept ----
\section{Proof-of-concept and on-chain evidence}\label{sec:poc}

\subsection{Foundry proof-of-concept}\label{sec:foundry}

Two independent Foundry test suites were written against unmodified contracts on Base mainnet forks. The first targets Slipstream CL gauges, the second PancakeSwap's MasterChefV3. Each suite conducts three tests proving the same three properties.

\subsubsection*{Slipstream CL gauges}

\textbf{Test 1 (no warmup).} A single-tick position was deposited into an active gauge and withdrawn after one block (two seconds). The position collected non-zero AERO rewards ($1.147 \times 10^{16}$ wei), confirming the absence of any warmup period.

\textbf{Test 2 (width independence).} A single-tick position and a ten-tick-spacing position were deposited simultaneously with equal capital. After ten blocks, the reward per unit of liquidity was identical for both positions ($2{,}527{,}437{,}520{,}521$ scaled reward units per $L$ for both, 0 basis points difference). The narrow position produced $20\times$ more liquidity per dollar, and therefore earned $20\times$ more AERO per dollar of capital deployed.

\textbf{Test 3 (duration independence).} A position staked for two blocks (4 seconds) and a position staked for twenty blocks (40 seconds) were compared. The reward rate per second per unit of liquidity was identical ($126{,}956{,}951{,}942$ scaled units per second per $L$ for both, 0 basis points difference), confirming that accrual rate is independent of staking duration.

\subsubsection*{PancakeSwap MasterChefV3}

The same three tests were reproduced against PancakeSwap's \texttt{MasterChefV3} (\texttt{0xC6A2...665A3}) on Base, using the WETH/USDC 500 fee tier pool (PID~5, allocPoint~678).

\textbf{Test 1 (no warmup).} A single-tick position was staked in MasterChefV3 via \texttt{safeTransferFrom} and harvested after one block (two seconds). The position received $2.707 \times 10^{15}$ CAKE wei, confirming non-zero rewards from a single block of staking.

\textbf{Test 2 (width independence).} A single-tick position and a twenty-tick position were minted with equal capital. After 60 seconds, the reward per unit of liquidity was identical for both (0 basis points difference: $662{,}204{,}883{,}480{,}521{,}107$ vs $662{,}204{,}883{,}480{,}520{,}996$ scaled reward units per $L$), confirming width independence. The narrow position produced $14\times$ more liquidity per dollar, and therefore earned $14\times$ more CAKE per dollar of capital deployed.

\textbf{Test 3 (duration independence).} A position staked for 4 seconds and a position staked for 40 seconds were compared. The reward rate per second per unit of liquidity was identical ($1{,}353{,}714{,}931{,}590{,}222$ CAKE wei per second per $L$ for both, 0 basis points difference), confirming no warmup period.

Both test suites pass against unmodified production contracts. The suites are available in the project repository (Appendix~\ref{app:foundry}) and are reproducible against any Base mainnet RPC endpoint.

\subsection{Aggregate on-chain evidence}\label{sec:onchain}

On-chain analysis of major CL pools across affected protocols, conducted between March and April 2026:

\begin{itemize}[leftmargin=2em]
\item In the largest WETH/USDC CL100 pool on one protocol, single-tick-width positions account for 63.4\% of in-range staked liquidity.
\item Across 50 sampled blocks on this pool, 46 (92\%) contained zero \texttt{Swap} events whilst single-tick positions remained staked. Emissions accrued in blocks where no swap was facilitated.
\item On a second protocol, single-tick positions in the four highest-emission pools capture an estimated 18.1\% of chain-level emission token distribution.
\item Across observed protocols and chains, positions matching the parasitic pattern are present in over 20 gauged pools simultaneously.
\item The strategy has been continuously active for months with no observed downtime.
\end{itemize}

All figures are derived from public on-chain data.

% ---- 5. Remediation ----
\section{Remediation}\label{sec:remediation}

\subsection{Ineffective approaches}\label{sec:ineffective}

\begin{table}[H]
\centering
\begin{tabular}{@{}ll@{}}
\toprule
Approach & Failure mode \\
\midrule
Address blacklisting & Contracts can be redeployed in one transaction \\
Global emission reduction & Penalises legitimate LPs equally \\
Manual intervention & Millions of autonomous transactions across continuous operation \\
\bottomrule
\end{tabular}
\end{table}

\subsection{Proposed gauge-level modifications}\label{sec:fixes}

\textbf{Fix 1: Minimum position width.} A width check at gauge deposit time:


\begin{center}\begin{minipage}{0.85\textwidth}\begin{verbatim}
require(tickUpper - tickLower >= minWidth, "Position too narrow");
\end{verbatim}\end{minipage}\end{center}


Fork testing against an active Slipstream gauge confirms that the capital efficiency penalty scales faster than linearly with enforced width. At \texttt{minWidth} = $2\times$ tick spacing, a parasitic operator needs $2\times$ more capital for the same liquidity. At $5\times$ tick spacing, the requirement rises to $6\times$. At $10\times$, $15\times$.

\textbf{Fix 2: Reward warmup period.} A linear ramp on reward accrual after each deposit:


\begin{center}\begin{minipage}{0.85\textwidth}\begin{verbatim}
uint256 elapsed = block.timestamp - depositTimestamp[tokenId];
uint256 scale = elapsed >= warmupPeriod
    ? 1e18
    : (elapsed * 1e18) / warmupPeriod;
earned = (rawEarned * scale) / 1e18;
\end{verbatim}\end{minipage}\end{center}


With a 60-second warmup, a position held for one block (2 seconds) earns at a scale factor of $2/60 = 3.3\%$ of its nominal rate. CL rebalancing requires minting a new NFT, so each rebalance restarts the warmup. The cost to legitimate LPs is the yield lost during the ramp. A linear ramp averages 50\% of nominal during the warmup window, so the drag is half the warmup divided by the rebalance interval: $30/1800 = 1.7\%$ for 30-minute rebalancing, $30/3600 = 0.8\%$ for hourly.

More comprehensive approaches incorporating position quality into the emission scoring function are possible but are beyond the scope of this paper.

\subsection{Combined effect}\label{sec:combined}

Fork testing both fixes together against an active Slipstream gauge produces the following results:

\begin{table}[H]
\centering
\begin{tabular}{@{}lr@{}}
\toprule
Scenario & Effective yield (\% of baseline) \\
\midrule
Parasitic (1 tick, 2s hold, no fixes) & 100\% \\
Parasitic ($5\times$ tick spacing min, 60s warmup) & 0.66\% \\
Legitimate LP (10 ticks, 30 min hold, 60s warmup) & 98.3\% \\
\bottomrule
\end{tabular}
\end{table}

The combined effect reduces parasitic extraction to less than 1\% of its unmitigated level while legitimate LPs retain over 98\% of their expected yield. The warmup alone shifts the yield ratio between parasitic and legitimate positions from 1:1 to 1:27. Neither fix in isolation is sufficient. Minimum width without warmup still permits rapid cycling. Warmup without width constraints still permits single-tick concentration.

% ---- 6. Affected Systems ----
\section{Affected systems}\label{sec:affected}

The vulnerability is present in any CL DEX that distributes emissions proportional to instantaneous staked liquidity without duration, width, or utilisation checks. The proof-of-concept suites in \S\ref{sec:foundry} verify this against two architecturally distinct implementations on Base. The capital requirement is a single position plus gas reserves, deployable from any address.

% ---- 7. Discussion ----
\section{Discussion}\label{sec:discussion}

\subsection{Parasitic liquidity as a distinct extraction class}\label{sec:distinct}

Parasitic liquidity differs from frontrunning, sandwich attacks, and other documented extraction techniques in three respects. First, it operates at the emission distribution layer rather than the swap execution layer. Second, it is invisible to standard DEX analytics. Pool TVL, volume, and fee metrics are unaffected. Only per-position emission breakdowns reveal the dilution. Third, it does not affect execution quality for traders. The cost falls entirely on legitimate LPs through reduced emission share.

Individual rebalance transactions are indistinguishable from legitimate LP adjustments. The aggregate pattern of per-block frequency, single-tick width, and zero-swap co-occurrence is required to identify the extraction.

\subsection{Implications for protocol design}\label{sec:implications}

Emission-bearing CL DEXs assume that directing emissions to a pool produces swap depth. Parasitic liquidity breaks this assumption. Whether emissions are allocated by voter direction (ve(3,3) models) or governance-set pool weights (MasterChef models), the intended incentive loop is the same:
$$\text{emission allocation} \xrightarrow{\text{rewards}} \text{LP} \xrightarrow{\text{depth}} \text{swaps} \xrightarrow{\text{fees}} \text{protocol/voters}$$
Parasitic liquidity interrupts this at the second step. Emissions flow to the position, but no depth is provided and no swap fees materialise. The emission becomes a subsidy to the operator, not a liquidity purchase.

\subsection{Allocation-layer optimisations}\label{sec:allocation}

Optimisations to how emissions are allocated across pools do not address how emissions are distributed within a pool. The vulnerability is in how rewards are split among positions within a pool, not in how pools are selected to receive them.

\subsection{Disclosure}\label{sec:disclosure}

The findings in this paper were disclosed to affected protocols via Immunefi prior to publication.

% ---- 8. Conclusion ----
\section{Conclusion}\label{sec:conclusion}

Parasitic liquidity is a structural vulnerability in the emission distribution layer of concentrated liquidity DEXs. Gauges measure instantaneous liquidity. Protocols intend to incentivise sustained swap depth. The gap between these two quantities is exploitable.

The vulnerability is demonstrated against unmodified mainnet contracts. Aggregate on-chain evidence confirms exploitation at scale across multiple protocols and chains. The strategy is linearly scalable with a low economic floor.

Remediation requires minimum position widths and reward warmup periods. Absent gauge-level modifications, replication is a matter of time.

% ---- References ----
\addcontentsline{toc}{section}{References}
\begin{thebibliography}{9}
\bibitem{adams2021}
Adams, H., Zinsmeister, N., Salem, M., Keefer, R. and Robinson, D. (2021) `Uniswap v3 Core', \textit{Uniswap Labs Technical Whitepaper}. Available at: \url{https://uniswap.org/whitepaper-v3.pdf}.
\end{thebibliography}

% ---- Appendices ----
\clearpage
\appendix

\section{Foundry test suites}\label{app:foundry}

The repository is at \url{https://github.com/g4nc0r/parasitic-liquidity} (archived at \url{https://doi.org/10.5281/zenodo.19528399}). Two independent Foundry projects fork Base mainnet and test against unmodified production contracts. Each requires a Base RPC endpoint via the \texttt{BASE\_RPC\_URL} environment variable.

\textbf{Suite 1: Slipstream CL gauges} (\texttt{slipstream/})

\begin{center}\begin{minipage}{0.85\textwidth}\begin{verbatim}
cd slipstream
BASE_RPC_URL=<your_base_rpc> forge test -vv
\end{verbatim}\end{minipage}\end{center}

\begin{itemize}[leftmargin=2em]
\item \texttt{SlipstreamParasitic.t.sol} --- no warmup, width independence, duration independence
\end{itemize}

\textbf{Suite 2: PancakeSwap MasterChefV3 + remediation} (\texttt{pcs/})

\begin{center}\begin{minipage}{0.85\textwidth}\begin{verbatim}
cd pcs
BASE_RPC_URL=<your_base_rpc> forge test -vv
\end{verbatim}\end{minipage}\end{center}

\begin{itemize}[leftmargin=2em]
\item \texttt{PCSParasitic.t.sol} --- no warmup, width independence, duration independence
\item \texttt{RemediationTest.t.sol} --- minimum width + warmup effectiveness (\S\ref{sec:combined})
\end{itemize}

\section{Data collection methodology}

Aggregate pool-level statistics were derived from:
\begin{itemize}[leftmargin=2em]
\item \texttt{cast call} (Foundry) for \texttt{stakedLiquidity}, \texttt{slot0}, and position parameters via NFPM \texttt{positions(tokenId)}
\item Block-level event logs for \texttt{Swap} event co-occurrence with gauge deposit and withdraw events
\item \texttt{balanceOf} on NFPM contracts to count positions held by high-frequency contracts
\item DexScreener for pool volume estimates
\end{itemize}

No individual addresses are reported. Pool-level aggregates were computed from public on-chain state.

\end{document}
